% Options for packages loaded elsewhere
\PassOptionsToPackage{unicode}{hyperref}
\PassOptionsToPackage{hyphens}{url}
\PassOptionsToPackage{dvipsnames,svgnames,x11names}{xcolor}
\documentclass[
  10pt,
  fontset=fandol]{ctexart}
\usepackage{xcolor}
\usepackage[margin=16mm]{geometry}
\usepackage{amsmath,amssymb}
\setcounter{secnumdepth}{-\maxdimen} % remove section numbering
\usepackage{iftex}
\ifPDFTeX
  \usepackage[T1]{fontenc}
  \usepackage[utf8]{inputenc}
  \usepackage{textcomp} % provide euro and other symbols
\else % if luatex or xetex
  \usepackage{unicode-math} % this also loads fontspec
  \defaultfontfeatures{Scale=MatchLowercase}
  \defaultfontfeatures[\rmfamily]{Ligatures=TeX,Scale=1}
\fi
\usepackage{lmodern}
\ifPDFTeX\else
  % xetex/luatex font selection
\fi
% Use upquote if available, for straight quotes in verbatim environments
\IfFileExists{upquote.sty}{\usepackage{upquote}}{}
\IfFileExists{microtype.sty}{% use microtype if available
  \usepackage[]{microtype}
  \UseMicrotypeSet[protrusion]{basicmath} % disable protrusion for tt fonts
}{}
\makeatletter
\@ifundefined{KOMAClassName}{% if non-KOMA class
  \IfFileExists{parskip.sty}{%
    \usepackage{parskip}
  }{% else
    \setlength{\parindent}{0pt}
    \setlength{\parskip}{6pt plus 2pt minus 1pt}}
}{% if KOMA class
  \KOMAoptions{parskip=half}}
\makeatother
\usepackage{longtable,booktabs,array}
\newcounter{none} % for unnumbered tables
\usepackage{calc} % for calculating minipage widths
% Correct order of tables after \paragraph or \subparagraph
\usepackage{etoolbox}
\makeatletter
\patchcmd\longtable{\par}{\if@noskipsec\mbox{}\fi\par}{}{}
\makeatother
% Allow footnotes in longtable head/foot
\IfFileExists{footnotehyper.sty}{\usepackage{footnotehyper}}{\usepackage{footnote}}
\makesavenoteenv{longtable}
\setlength{\emergencystretch}{3em} % prevent overfull lines
\providecommand{\tightlist}{%
  \setlength{\itemsep}{0pt}\setlength{\parskip}{0pt}}
\usepackage{amsmath,amssymb,mathtools,needspace}
\xeCJKDeclareCharClass{CJK}{"2032,"0400->"04FF,"2160->"216B,"2460->"2473,"25A0,"25C7,"25CB,"25CF}
\IfFontExistsTF{Arial Unicode MS}{\xeCJKsetup{AutoFallBack=true}\setCJKfallbackfamilyfont{\CJKrmdefault}{Arial Unicode MS}}{}
\setcounter{secnumdepth}{0}
\setlength{\emergencystretch}{2em}
\allowdisplaybreaks
\usepackage{bookmark}
\IfFileExists{xurl.sty}{\usepackage{xurl}}{} % add URL line breaks if available
\urlstyle{same}
\hypersetup{
  pdftitle={Taylor 级数法},
  colorlinks=true,
  linkcolor={Maroon},
  filecolor={Maroon},
  citecolor={Blue},
  urlcolor={Blue},
  pdfcreator={LaTeX via pandoc}}

\title{Taylor 级数法}
\author{}
\date{}

\begin{document}
\maketitle

\subsubsection{5.3.1 Taylor 级数法}\label{taylor-ux7ea7ux6570ux6cd5}

设初值问题 (5.1.1)、(5.1.2) 有解 \(y=y(x)\)，用 Taylor 展开，有

\[
y(x_{n+1})=y(x_n)+hy'(x_n)+\frac{h^2}{2!}y''(x_n)+\frac{h^3}{3!}y'''(x_n)+\cdots,\tag{5.3.1}
\]

其中 \(y(x)\) 的各阶导数依据所给方程 (5.1.1) 可以用函数 \(f\) 来表达。下面引进函数序列 \(f^{(j)}(x,y)\) 来描述求导过程，即

\[
y'=f\equiv f^{(0)},\qquad
y''=\frac{\partial f^{(0)}}{\partial x}+f\frac{\partial f^{(0)}}{\partial y}\equiv f^{(1)},\qquad
y'''=\frac{\partial f^{(1)}}{\partial x}+f\frac{\partial f^{(1)}}{\partial y}\equiv f^{(2)},\qquad\cdots.
\]

一般地

\[
y^{(j)}=\frac{\partial f^{(j-2)}}{\partial x}+f\frac{\partial f^{(j-2)}}{\partial y}\equiv f^{(j-1)},\qquad j=2,3,\cdots,
\]

而具体写出则有

\[
\begin{cases}
y'=f,\\
y''=\dfrac{\partial f}{\partial x}+f\dfrac{\partial f}{\partial y},\\
y'''=\dfrac{\partial^2f}{\partial x^2}+2f\dfrac{\partial^2f}{\partial x\partial y}
+f^2\dfrac{\partial^2f}{\partial y^2}
+\dfrac{\partial f}{\partial y}\left(\dfrac{\partial f}{\partial x}+f\dfrac{\partial f}{\partial y}\right),\\
\qquad\vdots
\end{cases}\tag{5.3.2}
\]

在展开式 (5.3.1) 的右端截取若干项，并且在 \((x_n,y_n)\) 按式 (5.3.2) 计算系数 \(y^{(j)}(x_n)\) 的近似值 \(y_n^{(j)}\)，结果导出下列 \textbf{Taylor 格式}

\[
y_{n+1}=y_n+hy'_n+\frac{h^2}{2!}y''_n+\cdots+\frac{h^p}{p!}y_n^{(p)},\tag{5.3.3}
\]

其中一阶 Taylor 格式 \((p=1)\)

\[
y_{n+1}=y_n+hy'_n
\]

就是 Euler 格式，而提高 Taylor 格式的阶 \(p\) 即可提高计算结果的精度。显然，\(p\) 阶

\href{../../source-images/pdf-127.jpeg}{原页}

Taylor 格式 (5.3.3) 的局部截断误差为

\[
y(x_{n+1})-y_{n+1}=\frac{h^{p+1}}{(p+1)!}y^{(p+1)}(\xi),\qquad x_n<\xi<x_{n+1},
\]

因此，它按下述定义具有 \(p\) 阶精度。

\textbf{定义 5.1} 如果一种方法的局部截断误差为 \(O(h^{p+1})\)，则称该方法具有 \(p\) 阶精度。

\textbf{例 5.3} 用 Taylor 级数法求解例 5.1 的初值问题 (5.2.2)。

\textbf{解} 直接求导知

\[
\begin{aligned}
y'&=y-\frac{2x}{y},\\
y''&=y'-\frac{2}{y^2}(y-xy'),\\
y'''&=y''+\frac{2}{y^2}(xy''+2y')-\frac{4x(y')^2}{y^3},\\
y^{(4)}&=y'''+\frac{2}{y^2}(xy'''+3y'')-\frac{12y'}{y^3}(xy''+y')+\frac{12x(y')^3}{y^4}.
\end{aligned}
\]

据此用四阶 Taylor 格式即可获得令人满意的结果（仍取步长 \(=0.1\) 计算，结果如表 5.3 所示）。

\textbf{表 5.3}

{\def\LTcaptype{none} % do not increment counter
\begin{longtable}[]{@{}lll@{}}
\toprule\noalign{}
\(x_n\) & \(y_n\) & \(y(x_n)\) \\
\midrule\noalign{}
\endhead
\bottomrule\noalign{}
\endlastfoot
0.1 & 1.0954 & 1.0954 \\
0.2 & 1.1832 & 1.1832 \\
0.3 & 1.2649 & 1.2649 \\
\end{longtable}
}

应当指出，Taylor 格式 (5.3.3) 从形式上看似简单，但具体构造这种格式往往是相当困难的，因为它需要按式 (5.3.2) 提供导数值 \(y_n^{(j)}\)。当阶数提高时，求导过程式 (5.3.2) 可能很复杂，因此 Taylor 级数法通常不直接使用，但可以用它来启发思路。

\end{document}
